\section{RQ4: 分析后训练中 CE Loss 的影响}

\subsection{RQ4 的比较对象}

这一节只讨论 fixed prefix-hint family 里的后训练 CE 系数，不把 dynamic CE trajectory 混进主表。对两个数据集，统一比较四条线：

\begin{itemize}
\item non-CE baseline：plain \texttt{fixed taskfix}
\item \texttt{CE=0.001}
\item \texttt{CE=0.005}
\item \texttt{CE=0.01}
\end{itemize}

因此 RQ4 回答的不是“加不加 CE”，而是“加了 CE 之后，系数怎么改变 top-10 / coverage 的形状，而且这种变化在 Instruments 与 Arts 上是否一致”。

\subsection{Instruments：0.005 提升 coverage，0.01 提升 long-run top-10}

\begin{table}[H]
\centering
\small
\setlength{\tabcolsep}{4pt}
\begin{tabular}{p{3.6cm} p{2.2cm} c c c c}
\toprule
Variant & Best ckpt & NDCG@10 & HR@10 & NDCG@50 & HR@50 \\
\midrule
fixed taskfix (no CE) & checkpoint-2997 & 0.0931 & 0.1189 & 0.1094 & 0.1941 \\
\texttt{CE=0.001} & checkpoint-2664 & 0.0931 & 0.1168 & 0.1102 & 0.1951 \\
\texttt{CE=0.005} & checkpoint-1665 & 0.0945 & 0.1180 & 0.1118 & 0.1985 \\
\texttt{CE=0.01} & checkpoint-2997 & 0.0953 & 0.1204 & 0.1113 & 0.1943 \\
\bottomrule
\end{tabular}
\caption{Instruments 上的 fixed CE coefficient sweep。}
\end{table}

\begin{figure}[H]
\centering
\includegraphics[width=0.96\textwidth]{ce_scaling_three_variants_curves.png}
\caption{Instruments 上 fixed CE coefficient 的完整曲线。}
\end{figure}

\begin{itemize}
\item \texttt{CE=0.001} 更像一个温和正则：\texttt{NDCG@10} 基本不变，但 \texttt{HR@50} 从 \texttt{0.1941} 提到 \texttt{0.1951}。
\item \texttt{CE=0.005} 是当前 coverage 最强点，\texttt{HR@50} 来到 \texttt{0.1985}，同时 \texttt{NDCG@10} 也升到 \texttt{0.0945}。
\item \texttt{CE=0.01} 则把 best \texttt{NDCG@10} 继续抬到 \texttt{0.0953}，而且 best 点回到更晚的 \texttt{checkpoint-2997}；它更像“long-run top-10 强化”，而不是 coverage 最优。
\end{itemize}

\subsection{Arts：0.001 的 coverage 最强，但 0.005 更稳地保住 headline}

\begin{table}[H]
\centering
\small
\setlength{\tabcolsep}{4pt}
\begin{tabular}{p{3.6cm} p{2.2cm} c c c c}
\toprule
Variant & Best ckpt & NDCG@10 & HR@10 & NDCG@50 & HR@50 \\
\midrule
fixed taskfix (no CE) & checkpoint-3368 & 0.0956 & 0.1248 & 0.1120 & 0.2007 \\
\texttt{CE=0.001} & checkpoint-2105 & 0.0943 & 0.1244 & 0.1114 & 0.2029 \\
\texttt{CE=0.005} & checkpoint-2105 & 0.0956 & 0.1257 & 0.1118 & 0.2006 \\
\texttt{CE=0.01} & checkpoint-2105 & 0.0941 & 0.1228 & 0.1106 & 0.1990 \\
\bottomrule
\end{tabular}
\caption{Arts 上的 fixed CE coefficient sweep。}
\end{table}

\begin{figure}[H]
\centering
\includegraphics[width=0.96\textwidth]{genrec-only-arts-fixed-ce-coefficients-curves.png}
\caption{Arts 上 fixed CE coefficient 的完整曲线。}
\end{figure}

\begin{itemize}
\item Arts 上最明显的 coverage 增益来自 \texttt{CE=0.001}：\texttt{HR@50} 从 \texttt{0.2007} 提到 \texttt{0.2029}，但 top-10 略有回落。
\item \texttt{CE=0.005} 则保住了和 non-CE baseline 持平的 \texttt{NDCG@10=0.0956}，同时把 \texttt{HR@10} 拉到四条线里最高的 \texttt{0.1257}。
\item \texttt{CE=0.01} 在 Arts 上反而没有复制 Instruments 的 long-run top-10 优势，因此它更像一个数据集敏感的高倍率设定。
\end{itemize}

\subsection{跨数据集读法}

\begin{table}[H]
\centering
\small
\begin{tabular}{p{2.2cm} p{3.0cm} p{3.0cm} p{5.0cm}}
\toprule
Dataset & best top-10 coefficient & best coverage coefficient & 读法 \\
\midrule
Instruments & \texttt{CE=0.01} & \texttt{CE=0.005} & 更高倍率会把 top-10 往长跑尾段抬，但 coverage 峰值仍出现在中等系数 \\
Arts & \texttt{CE=0.005} 与 no-CE 并列 & \texttt{CE=0.001} & 较小系数更像 coverage regularizer；过大的 \texttt{0.01} 会同时伤到 top-10 和 coverage \\
\bottomrule
\end{tabular}
\caption{CE 系数的跨数据集 summary。}
\end{table}

因此，RQ4 的主结论不是“存在一个在所有数据集都最优的 CE 系数”，而是：

\begin{itemize}
\item CE 确实会系统性改变 fixed prefix-hint family 的后训练轨迹；
\item 但它改变的是 top-10 与 coverage 的形状，而不是单调地“系数越大越好”；
\item 当前最稳妥的说法是：\texttt{0.005} 更像跨数据集都还站得住的中间档，\texttt{0.01} 更像 Instruments 特有的 top-10 aggressive setting，\texttt{0.001} 则更像 Arts 上的 mild coverage regularizer。
\end{itemize}
